%!TEX root = ../main.tex

\section{Results}
\label{sec:results}

This section reports the empirical evaluation of the proposed trust layer --- a
split-conformal heading-interval and risk-controlled steering policy bolted onto
a \emph{stock} detector (YOLOv8n) with a robust Huber-IRLS line fit. \emph{No
dataset-level numbers are filled in yet}: every quantitative cell is a
placeholder (``\texttt{--}'' in tables, ``[TODO]'' in prose) to be completed once
the pre-registered runs (experiments E0--E6) have finished. The table
structures, comparison arms, metrics, statistical tests, and pass criteria below
are fixed in advance so that the analysis is confirmatory rather than
exploratory. Each table and figure is wrapped in \verb|\IfFileExists| so the
manuscript compiles today; the auto-generated file is \verb|\input| once a run
produces it.

We restate the readout so the tables are self-contained. From the post-NMS box
centroids $\{(x_i,y_i)\}_{i=1}^{n}$ we fit the guidance line $x = a\,y + b$ by
robust Huber-IRLS, with weights $w_i$ from the Huber influence function; the
heading is $\theta = \arctan(a)$ in degrees from vertical and the cross-track
offset is $e_{ct}$ (px, or cm only under the calibrated homography of
Section~\ref{sec:setup-homography}). This readout adds no NPU operators because
it is CPU math executed after NMS; the exported detector graph is a vanilla YOLO.
The trust score is the leave-one-out (LOO) heading dispersion $s$ (deg): the
dispersion of $\theta$ over refits that each omit one centroid, which targets
wrong-row contamination without any extra labels. Split-conformal calibration on
a held-out score set $\{s_1,\dots,s_m\}$ (split \emph{by base image}) gives the
quantile $\hat q$, and the heading interval is $[\theta - h, \theta + h]$ with
half-width $h = \hat q\,s$, locally adaptive in the sense of conformalized
quantile regression~\cite{romano2019cqr,lei2018,vovk2005,angelopoulos2023}. We
report nominal miscoverage $\alpha\in\{0.10,0.05\}$ and the realised coverage,
with Clopper--Pearson / Wilson intervals. Throughout, the detector's native DFL
box-edge distribution variance $\sigma^2_{\mathrm{DFL}}$~\cite{gfl,gflv2} appears
\emph{only} as a non-informative negative control; the deployed line fit is
robust Huber-IRLS, not variance weighting.

\paragraph{Scope and honesty caveats.} Three caveats bound every claim below.
(i)~This is an \emph{applied} contribution --- a trust layer and evaluation on a
stock detector --- not a new detector architecture. (ii)~The conformal coverage
guarantee is finite-sample-valid only \emph{under exchangeability within one
distribution}; results on CRBD and on field video are therefore a
coverage-\emph{degradation} study, not a guarantee, and are labelled as such
wherever coverage is reported. (iii)~The policy replay (E6) is an
\emph{open-loop} kinematic-bicycle replay~\cite{rajamani2012}, not a closed-loop
field-robot run: no odometry is in the loop, so E6 measures how the abstain
policy would have filtered large-heading-error frames, not realised tracking
error. Centimetre units are used only where the camera is calibrated (the robot
platform and the CRDLD homography); field video stays in px and deg.

\subsection{E0 --- Is the trust score informative? (and the DFL-variance negative control)}
\label{sec:res-e0-trustinfo}

E0 tests the premise of the whole trust layer: that the LOO dispersion score $s$
actually tracks the realised heading error, frame by frame. Table~\ref{tab:e0-trustinfo}
reports the Spearman rank correlation $\rho(s,\,|\theta-\theta^\star|)$ between
the score and the realised heading error on CRDLD test (split by base image,
$\ge 5$ seeds), with a bootstrap CI. The pass criterion is $\rho > 0$ with a CI
excluding zero: a positive, significant correlation is what licenses using $s$
as a nonconformity score in E2--E4. The same table reports, as a \emph{negative
control}, the rank correlation of the detector's native DFL box-edge variance
$\sigma^2_{\mathrm{DFL}}$~\cite{gfl,gflv2} with the identical error target; our
prior measurement on this data is $\rho_{\mathrm{DFL}} = -0.13$, i.e.\
non-informative (and, if anything, anti-correlated). This contrast is the
empirical reason the deployed readout uses robust Huber-IRLS rather than
DFL-variance precision weighting: the box-distribution variance does not track
heading reliability here. Figure~\ref{fig:e0-scatter} visualises both
relationships as score-versus-error scatters.

\IfFileExists{tables/tab_e0.tex}{\input{tables/tab_e0.tex}}{%
{\renewcommand{\arraystretch}{1.3}
\begin{table}[!ht]
  \centering
  \caption{E0 --- Trust-score informativeness on CRDLD test, split by base
    image, $\ge 5$ seeds. Spearman $\rho$ between each candidate score and the
    realised heading error $|\theta-\theta^\star|$, with bootstrap 95\,\% CIs.
    The proposed LOO-dispersion score $s$ must satisfy $\rho>0$ (CI excluding
    zero); the DFL-variance row $\sigma^2_{\mathrm{DFL}}$ is the non-informative
    negative control. \emph{Results to be completed.}}
  \label{tab:e0-trustinfo}
  \resizebox{\columnwidth}{!}{\begin{tabular}{lcccc}
    \hline
    \textbf{Candidate score} & \textbf{Spearman $\rho$} & \textbf{95\,\% CI} & \textbf{$\rho>0$?} & \textbf{Role} \\
    \hline
    $s$ = LOO heading dispersion        & -- & [--,\,--] & -- & proposed \\
    Fit residual (median)               & -- & [--,\,--] & -- & alt.\ score \\
    Detector confidence$^2$             & -- & [--,\,--] & -- & alt.\ score \\
    RANSAC consensus fraction           & -- & [--,\,--] & -- & alt.\ score \\
    $\sigma^2_{\mathrm{DFL}}$ (DFL var) & -- & [--,\,--] & -- & negative control \\
    \hline
  \end{tabular}}
  \medskip
\end{table}}%
}

\IfFileExists{figures/fig_e0_scatter.pdf}{%
\begin{figure}[!ht]\centering
\includegraphics[width=0.72\linewidth]{figures/fig_e0_scatter.pdf}
\caption{E0 --- Realised heading error versus candidate trust score on CRDLD
  test: the proposed LOO-dispersion score $s$ (informative) contrasted with the
  DFL box-edge variance $\sigma^2_{\mathrm{DFL}}$ (non-informative negative
  control). \emph{Auto-generated from the run.}}\label{fig:e0-scatter}
\end{figure}}{}

\subsection{E1 --- Point-accuracy floor of the guidance readout}
\label{sec:res-e1-accuracy}

E1 establishes the point-heading accuracy floor: before any interval is claimed,
the deployed Huber-IRLS readout must produce a competitive point heading.
Table~\ref{tab:e1-accuracy} reports heading error (deg) on CRDLD test ($\ge 5$
seeds), and line-fit RMSE (px, and cm under the calibrated homography), for six
readout arms that all share a single stock YOLOv8n checkpoint and differ only in
the CPU-side line fit: equal-weight least squares (\textbf{equalLS}, the
Diao-style detect-then-fit recipe~\cite{diao_yolox}), \textbf{RANSAC} least
squares, the deployed robust \textbf{IRLS} (Huber), the dead DFL-variance
precision-weighted variant (\textbf{CALM}, included only as the negative control
of E0/E3), a qualitative single-frame heuristic (\textbf{qual}), and an
\textbf{oracle} that picks the best-of-arm per frame (an upper bound, not a
deployable method). Because the arms are paired per frame, comparisons use the
Wilcoxon signed-rank test with Holm correction across the confirmatory family
and Cliff's $\delta$ as effect size. The pre-registered claims are: IRLS
\emph{matches} RANSAC (no significant difference, or $|\delta|<0.147$) and IRLS
\emph{far beats} equalLS (Holm-adjusted $p<0.05$ and $|\delta|\ge 0.33$), with
median heading sub-degree. We also report cross-dataset transfer (CRDLD
$\leftrightarrow$ CRBD~\cite{crbd,cth1,cth2}) so the floor is not a single-domain
artefact. Figure~\ref{fig:e1-box} shows the per-arm heading-error distribution,
and Figure~\ref{fig:qual-lines} shows qualitative fits.

\IfFileExists{tables/tab_e1.tex}{\input{tables/tab_e1.tex}}{%
{\renewcommand{\arraystretch}{1.3}
\begin{table}[!ht]
  \centering
  \caption{E1 --- Point-accuracy floor on CRDLD test, mean$\pm$std over $\ge 5$
    seeds; cm only under the calibrated homography (Section~\ref{sec:setup-homography}).
    All arms share one stock YOLOv8n checkpoint and differ only in the post-NMS
    line fit. Pairwise tests vs.\ the deployed IRLS arm: Wilcoxon signed-rank,
    Holm-adjusted, Cliff's $\delta$. \texttt{calm}=DFL-variance weighting
    (negative control); \texttt{oracle}=best-of-arm upper bound (not
    deployable). \emph{Results to be completed.}}
  \label{tab:e1-accuracy}
  \resizebox{\columnwidth}{!}{\begin{tabular}{lccccc}
    \hline
    \textbf{Readout arm} & \textbf{Heading ($^\circ$)$\downarrow$} & \textbf{RMSE (px)$\downarrow$} & \textbf{RMSE (cm)$\downarrow$} & \textbf{Wilcoxon $p$ (Holm)} & \textbf{Cliff's $\delta$} \\
    \hline
    equalLS~\cite{diao_yolox}        & -- & -- & -- & -- & -- \\
    RANSAC                           & -- & -- & -- & -- & -- \\
    \textbf{IRLS (deployed, Huber)}  & -- & -- & -- & \emph{ref} & \emph{ref} \\
    calm (DFL-var, neg.\ control)    & -- & -- & -- & -- & -- \\
    qual (single-frame heuristic)    & -- & -- & -- & -- & -- \\
    oracle (best-of-arm, upper bound)& -- & -- & -- & -- & -- \\
    \hline
  \end{tabular}}
  \medskip
\end{table}}%
}

\IfFileExists{figures/fig_e1_box.pdf}{%
\begin{figure}[!ht]\centering
\includegraphics[width=0.62\linewidth]{figures/fig_e1_box.pdf}
\caption{E1 --- Per-arm heading-error distribution on CRDLD test ($\ge 5$
  seeds). Boxes show median and IQR; the deployed IRLS arm is highlighted.
  \emph{Auto-generated from the run.}}\label{fig:e1-box}
\end{figure}}{}

\IfFileExists{figures/fig_qualitative.pdf}{%
\begin{figure}[!ht]\centering
\includegraphics[width=0.92\linewidth]{figures/fig_qualitative.pdf}
\caption{Qualitative guidance lines on CRDLD test frames: ground-truth line
  (green) vs.\ the deployed Huber-IRLS fit (red) over detected row-segment
  centroids, with the conformal heading interval shaded.
  \emph{Auto-generated from the run.}}\label{fig:qual-lines}
\end{figure}}{}

\subsection{E2 --- Realised vs.\ nominal coverage}
\label{sec:res-e2-coverage}

E2 is the central validity check for the trust layer. Table~\ref{tab:e2-coverage}
reports the realised coverage --- the empirical fraction of test frames whose
true heading lies in $[\theta-h,\theta+h]$ --- against the nominal target
$1-\alpha$ for $\alpha\in\{0.10,0.05\}$, with Clopper--Pearson and Wilson 95\,\%
CIs. Calibration and test are split \emph{by base image} so that exchangeability
is not broken by near-duplicate frames; this split is what makes the
split-conformal coverage guarantee~\cite{vovk2005,lei2018,angelopoulos2023}
applicable in-domain. The pass criterion is that the realised-coverage CI
contains the nominal target on the CRDLD in-domain test (marginal coverage), and
we additionally report coverage stratified by difficulty (near/far rows, low/high
centroid count) to expose conditional under-coverage. Cross-dataset rows
(CRDLD$\to$CRBD and field video~\cite{crbd,cth2}) are reported as a
coverage-\emph{degradation} study: exchangeability does not hold across
distributions, so any shortfall there quantifies how far the in-domain guarantee
travels, and is \emph{not} a claimed guarantee. Figure~\ref{fig:e2-reliability}
is the reliability diagram (realised vs.\ nominal across a sweep of $\alpha$).
We position this as a domain instantiation of conformal prediction on a detector
and explicitly contrast, rather than claim primacy over, prior
conformal-on-detector work~\cite{confdet1,confdet2}.

\IfFileExists{tables/tab_e2.tex}{\input{tables/tab_e2.tex}}{%
{\renewcommand{\arraystretch}{1.3}
\begin{table}[!ht]
  \centering
  \caption{E2 --- Realised vs.\ nominal heading-interval coverage, split by base
    image, with Clopper--Pearson / Wilson 95\,\% CIs. CRDLD is the in-domain
    guarantee (CI should contain $1-\alpha$); CRDLD$\to$CRBD and field video are
    a coverage-\emph{degradation} study under distribution shift (not a
    guarantee). \emph{Results to be completed.}}
  \label{tab:e2-coverage}
  \resizebox{\columnwidth}{!}{\begin{tabular}{llcccc}
    \hline
    \textbf{Eval set} & \textbf{$\alpha$} & \textbf{Nominal $1-\alpha$} & \textbf{Realised} & \textbf{Clopper--Pearson CI} & \textbf{In CI?} \\
    \hline
    \multirow{2}{*}{CRDLD (in-domain)}
      & 0.10 & 0.90 & -- & [--,\,--] & -- \\
      & 0.05 & 0.95 & -- & [--,\,--] & -- \\
    \hline
    \multirow{2}{*}{CRDLD$\to$CRBD (shift)}
      & 0.10 & 0.90 & -- & [--,\,--] & -- \\
      & 0.05 & 0.95 & -- & [--,\,--] & -- \\
    \hline
    \multirow{2}{*}{Field video (shift)}
      & 0.10 & 0.90 & -- & [--,\,--] & -- \\
      & 0.05 & 0.95 & -- & [--,\,--] & -- \\
    \hline
  \end{tabular}}
  \medskip
\end{table}}%
}

\IfFileExists{figures/fig_e2_reliability.pdf}{%
\begin{figure}[!ht]\centering
\includegraphics[width=0.55\linewidth]{figures/fig_e2_reliability.pdf}
\caption{E2 --- Reliability diagram: realised vs.\ nominal heading-interval
  coverage over a sweep of $\alpha$, CRDLD in-domain (on the diagonal under
  exchangeability) with the distribution-shift sets shown as a degradation
  study. \emph{Auto-generated from the run.}}\label{fig:e2-reliability}
\end{figure}}{}

\subsection{E3 --- Sharpness: which score gives the tightest valid interval?}
\label{sec:res-e3-sharpness}

A trust layer is useful only if its valid intervals are also \emph{sharp}: wide
intervals that always cover are uninformative for steering. E3 ablates the
nonconformity score, holding the conformal machinery and the target coverage
fixed, and reports the mean interval half-width $h$ at matched realised coverage.
Table~\ref{tab:e3-sharpness} compares the proposed LOO-dispersion score $s$
against the fit residual, the squared detector confidence, the DFL box-edge
variance $\sigma^2_{\mathrm{DFL}}$~\cite{gfl,gflv2}, and the RANSAC consensus
fraction. Because all scores are conformalized to the same $1-\alpha$, coverage
is (by construction) matched and the discriminating axis is width: the
pre-registered claim is that the failure-mode-matched LOO-dispersion score yields
the \emph{narrowest} mean half-width at matched coverage, with the
non-informative DFL-variance score giving the widest. Figure~\ref{fig:e3-width}
plots the half-width distribution per score.

\IfFileExists{tables/tab_e3.tex}{\input{tables/tab_e3.tex}}{%
{\renewcommand{\arraystretch}{1.3}
\begin{table}[!ht]
  \centering
  \caption{E3 --- Interval sharpness (score ablation) on CRDLD test, split by
    base image, $\ge 5$ seeds, at matched realised coverage ($\alpha=0.10$).
    Mean conformal half-width $h$ (deg, lower is sharper); realised coverage
    column confirms the comparison is at matched coverage. The proposed
    LOO-dispersion score $s$ should be tightest; $\sigma^2_{\mathrm{DFL}}$
    widest. \emph{Results to be completed.}}
  \label{tab:e3-sharpness}
  \resizebox{\columnwidth}{!}{\begin{tabular}{lccc}
    \hline
    \textbf{Nonconformity score} & \textbf{Realised coverage} & \textbf{Mean half-width $h$ ($^\circ$)$\downarrow$} & \textbf{Rank} \\
    \hline
    $s$ = LOO heading dispersion        & -- & -- & -- \\
    Fit residual                        & -- & -- & -- \\
    Detector confidence$^2$             & -- & -- & -- \\
    $\sigma^2_{\mathrm{DFL}}$ (DFL var) & -- & -- & -- \\
    RANSAC consensus fraction           & -- & -- & -- \\
    \hline
  \end{tabular}}
  \medskip
\end{table}}%
}

\IfFileExists{figures/fig_e3_width.pdf}{%
\begin{figure}[!ht]\centering
\includegraphics[width=0.62\linewidth]{figures/fig_e3_width.pdf}
\caption{E3 --- Conformal half-width distribution per nonconformity score at
  matched coverage on CRDLD test; narrower is sharper.
  \emph{Auto-generated from the run.}}\label{fig:e3-width}
\end{figure}}{}

\subsection{E4 --- Selective risk: does the conformal gate beat a confidence gate?}
\label{sec:res-e4-selrisk}

E4 evaluates the steering policy as a selective predictor~\cite{geifman2017selective,angelopoulos2022rcps}:
frames are accepted (proceed) or aborted (abstain), and we plot the heading error
on \emph{accepted} frames against the abort rate. Table~\ref{tab:e4-selrisk}
compares the conformal-width gate (abstain when the half-width $h$ exceeds a
threshold) against the fixed-confidence baseline gate (abstain when detector
confidence is below a threshold), at \emph{matched abort rate} so the comparison
is fair. The pre-registered claim is that, at matched abort, the conformal-width
gate yields lower heading error on accepted frames --- i.e.\ it aborts the right
frames --- because it is driven by a failure-mode-matched score rather than by
raw confidence. Figure~\ref{fig:e4-riskcov} is the risk--coverage curve (accepted
error vs.\ abort rate) for both gates; a uniformly lower conformal curve is the
visual statement of the same claim.

\IfFileExists{tables/tab_e4.tex}{\input{tables/tab_e4.tex}}{%
{\renewcommand{\arraystretch}{1.3}
\begin{table}[!ht]
  \centering
  \caption{E4 --- Selective risk on CRDLD test, split by base image, $\ge 5$
    seeds. Heading error on accepted frames at matched abort rate: conformal-width
    gate vs.\ fixed-confidence gate. Lower accepted error at equal abort is
    better. \emph{Results to be completed.}}
  \label{tab:e4-selrisk}
  \resizebox{\columnwidth}{!}{\begin{tabular}{lcccc}
    \hline
    \textbf{Abort rate} & \textbf{Conf-gate error ($^\circ$)$\downarrow$} & \textbf{Conformal-gate error ($^\circ$)$\downarrow$} & \textbf{$\Delta$ ($^\circ$)} & \textbf{Better} \\
    \hline
    0\,\% (no abstain) & -- & -- & -- & -- \\
    5\,\%              & -- & -- & -- & -- \\
    10\,\%             & -- & -- & -- & -- \\
    20\,\%             & -- & -- & -- & -- \\
    \hline
  \end{tabular}}
  \medskip
\end{table}}%
}

\IfFileExists{figures/fig_e4_riskcoverage.pdf}{%
\begin{figure}[!ht]\centering
\includegraphics[width=0.62\linewidth]{figures/fig_e4_riskcoverage.pdf}
\caption{E4 --- Selective risk: heading error on accepted frames vs.\ abort
  rate, conformal-width gate vs.\ fixed-confidence gate at matched abort.
  \emph{Auto-generated from the run.}}\label{fig:e4-riskcov}
\end{figure}}{}

\subsection{E5 --- Edge deployment: latency, throughput, and coverage under INT8}
\label{sec:res-e5-deploy}

E5 supports the deployment argument. Table~\ref{tab:e5-deploy} reports latency
(mean$\pm$std and P95), throughput (FPS), and energy per frame on the two named
edge stacks --- TensorRT (FP16/INT8) on the Jetson Nano~\cite{tensorrt} and RKNN
INT8 on the Rockchip board~\cite{rknntoolkit} --- together with the post-NMS CPU
overhead of the trust layer measured separately. Because the readout and trust
layer add no NPU operators (CPU, post-NMS), the exported INT8 graph is a vanilla
YOLO and the NPU latency is unchanged by the trust layer; the table makes this
verifiable by reporting the CPU overhead as a distinct row. The table's last
columns report the realised coverage \emph{after} INT8 post-training
quantization~\cite{jacob2018int8} --- recomputing the conformal quantile $\hat q$
on the quantized model's calibration scores --- to demonstrate that the coverage
guarantee survives quantization (the realised-coverage CI should still contain
$1-\alpha$ in-domain after the re-fit). This is a re-fit-after-INT8 result, not a
transfer of the FP32 quantile.

\IfFileExists{tables/tab_e5.tex}{\input{tables/tab_e5.tex}}{%
{\renewcommand{\arraystretch}{1.3}
\begin{table}[!ht]
  \centering
  \caption{E5 --- Edge latency, throughput, energy, and coverage-after-INT8.
    TensorRT on Jetson Nano; RKNN INT8 on Rockchip. Input at the stated
    resolution, batch size 1, after warm-up. The trust layer adds no NPU
    operators; its cost is the post-NMS CPU row. Coverage-after-INT8 re-fits
    $\hat q$ on the quantized model (CI should contain $1-\alpha$ in-domain).
    \emph{Results to be completed.}}
  \label{tab:e5-deploy}
  \resizebox{\columnwidth}{!}{\begin{tabular}{llccccc}
    \hline
    \textbf{Device / runtime} & \textbf{Precision} & \textbf{Latency mean$\pm$std (ms)$\downarrow$} & \textbf{P95 (ms)} & \textbf{FPS$\uparrow$} & \textbf{Energy/frame (mJ)$\downarrow$} & \textbf{Coverage @$\alpha{=}0.10$} \\
    \hline
    Jetson Nano / TensorRT & FP16 & -- & -- & -- & -- & -- \\
    Jetson Nano / TensorRT & INT8 & -- & -- & -- & -- & -- \\
    Rockchip / RKNN        & INT8 & -- & -- & -- & -- & -- \\
    \hline
    \multicolumn{2}{l}{Trust layer (post-NMS CPU overhead)} & -- & -- & -- & -- & n/a \\
    \hline
  \end{tabular}}
  \medskip
\end{table}}%
}

\subsection{E6 --- Open-loop policy replay}
\label{sec:res-e6-replay}

E6 asks whether the abstain policy would have improved guidance in a sequential
setting. \emph{This is an open-loop replay, not a closed-loop field run}: there
is no odometry and no real robot in the loop. We replay a held-out video sequence
through a kinematic bicycle model~\cite{rajamani2012}, feeding the per-frame
heading either unconditionally (fixed-threshold baseline) or through the
risk-controlled policy (proceed if $h\le\tau_p$; slow if $\tau_p<h\le\tau_s$;
abstain if $h>\tau_s$), and we count how many large-heading-error frames are
accepted under each policy. Table~\ref{tab:e6-replay} reports the fraction of
accepted frames whose heading error exceeds a safety threshold, plus the abstain
and slow rates; the pre-registered claim is that the conformal abstain policy
reduces the rate of accepted large-error frames relative to the fixed threshold
at a comparable abstain budget. Figure~\ref{fig:e6-replay} overlays the replayed
heading trace, the conformal interval, and the policy state over time. We
emphasise this is a replay-based plausibility result; a closed-loop field-robot
evaluation is left to future work and is stated as a limitation.

\IfFileExists{tables/tab_e6.tex}{\input{tables/tab_e6.tex}}{%
{\renewcommand{\arraystretch}{1.3}
\begin{table}[!ht]
  \centering
  \caption{E6 --- Open-loop kinematic-bicycle replay (\emph{not} closed-loop; no
    odometry) on a held-out field-video sequence. Fraction of accepted frames
    with heading error above the safety threshold, plus abstain/slow rates:
    fixed-threshold baseline vs.\ the risk-controlled conformal policy at a
    comparable abstain budget. Lower accepted large-error rate is better.
    \emph{Results to be completed.}}
  \label{tab:e6-replay}
  \resizebox{\columnwidth}{!}{\begin{tabular}{lcccc}
    \hline
    \textbf{Policy} & \textbf{Accepted large-error rate$\downarrow$} & \textbf{Abstain rate} & \textbf{Slow rate} & \textbf{Mean accepted error ($^\circ$)$\downarrow$} \\
    \hline
    Fixed-threshold baseline          & -- & -- & n/a & -- \\
    Risk-controlled conformal policy  & -- & -- & --  & -- \\
    \hline
  \end{tabular}}
  \medskip
\end{table}}%
}

\IfFileExists{figures/fig_e6_replay.pdf}{%
\begin{figure}[!ht]\centering
\includegraphics[width=0.78\linewidth]{figures/fig_e6_replay.pdf}
\caption{E6 --- Open-loop replay trace: per-frame heading with its conformal
  interval and the policy state (proceed/slow/abstain) over time, contrasted with
  the fixed-threshold baseline. Replay only; no odometry in the loop.
  \emph{Auto-generated from the run.}}\label{fig:e6-replay}
\end{figure}}{}

\subsection{Dataset and summary of planned findings}
\label{sec:res-summary}

Figure~\ref{fig:dataset} shows representative frames from the public field-video
dataset (real maize video) with its hand-labelled evaluation subset (C3), which
underpins the field-video rows of E2 and the sequence in E6. Taken together, the
experiments answer the questions the trust layer must pass before deployment.
\textbf{E0} establishes that the LOO-dispersion score is informative
($\rho>0$) and that the DFL box-edge variance is not (the
$\rho_{\mathrm{DFL}}=-0.13$ negative control), which is the empirical reason the
readout is robust Huber-IRLS rather than variance weighting. \textbf{E1} fixes
the point-accuracy floor (IRLS matches RANSAC, far beats equalLS, sub-degree
median). \textbf{E2} checks that conformal coverage holds in-domain under
exchangeability and quantifies its degradation under shift. \textbf{E3} shows
the proposed score gives the sharpest valid interval. \textbf{E4} shows the
conformal gate beats a fixed-confidence gate at matched abort. \textbf{E5}
reports edge latency/throughput and that coverage survives INT8 re-fit.
\textbf{E6} shows, in open-loop replay, that the abstain policy filters
large-error frames. Every cell will be populated once the pre-registered runs
complete; no dataset-level number is asserted here.

\IfFileExists{figures/fig_dataset.pdf}{%
\begin{figure}[!ht]\centering
\includegraphics[width=0.92\linewidth]{figures/fig_dataset.pdf}
\caption{Public crop-row field-video dataset (real maize video) with the
  hand-labelled evaluation subset used in E2 (field-video coverage) and E6
  (replay sequence). \emph{Auto-generated from the run.}}\label{fig:dataset}
\end{figure}}{}
